\section{相关工作}
\label{sec:related}

\subsection{面向 coding agent 的 harness 工程与评测}
\label{ssec:related_harness}
harness 工程指的是设计模型周边系统的实践，包括其工具、接口、记忆、执行约束与反馈闭环，它们共同塑造了智能体在长程任务上能够做什么~\citep{rajasekaranHarnessDesignLongrunning2026,lopopoloHarnessEngineeringLeveraging2026,trivedyImprovingDeepAgents2026,anthropicClaudecode2025,steinbergerOpenClawPersonalAI2026,HermesAgentAgent}。
具体而言，harness 中介了模型如何感知环境并在其上行动：它暴露出工具增强推理所依托的动作与观测接口~\citep{anthropicClaudecode2025}、用于仓库导航、文件编辑与命令执行的定制智能体—计算机接口~\citep{yangSWEagentAgentComputerInterfaces2024}，以及使长程运行保持可复现的沙箱化执行与编排支持~\citep{wangOpenHandsOpenPlatform2025}。

为验证这类系统是否真的有帮助，coding agent 的评测沿两条轴线并行走向成熟：任务时程与环境真实性。
其覆盖范围从关注污染与时效性控制的短程函数级基准~\citep{zhuoBigCodeBenchBenchmarkingCode2024,jainLiveCodeBenchHolisticContamination2024}，经由仓库规模的可执行补丁求解~\citep{jimenezSWEbenchCanLanguage2023,yangSWEbenchMultimodalAI2024,dengSWEBenchProCan2025}，一直延伸到考察长程、真实执行的多小时终端驱动工作流~\citep{miserendinoSWELancerCanFrontier2025,chanMLEbenchEvaluatingMachine2024,merrillTerminalBenchBenchmarkingAgents2026a}。
与之并行的基础设施路线围绕这些基准封装了可执行运行时与验证器~\citep{panTrainingSoftwareEngineering2025,jainR2EGymProceduralEnvironment2025,zengSWEHubUnifiedProduction2026}，其对可复现、可追溯、可验证执行的重视，正是 AHE 所依托的观测系统的直接动机。


\subsection{LLM 智能体的自动化优化}
\label{ssec:related_auto}

自动化智能体优化的各种方法，差别在于优化器观测到什么证据、以及它能够编辑什么。
一些工作通过按回合（episodic）的批评与反思来修订智能体自身的输出~\citep{madaanSelfRefineIterativeRefinement2023c,shinnReflexionLanguageAgents2023c,guoCritiQMiningData2025}。
另一些工作则以提示词与指令为目标~\citep{khattabDSPyCompilingDeclarative2023}：结构化的操作手册~\citep{zhangAgenticContextEngineering2025a}、语义优势先验~\citep{caiTrainingFreeGroupRelative2025}、面向多阶段程序的指令与示例联合优化流水线~\citep{opsahl-ongOptimizingInstructionsDemonstrations2024a}，以及由帕累托前沿轨迹驱动的反思式更新~\citep{agrawalGEPAReflectivePrompt2025a}。
还有一条独立的路线直接编辑程序结构本身，其形式包括技能库~\citep{wangVoyagerOpenEndedEmbodied2023a}、通过变异演化并打分的程序与智能体档案库~\citep{novikovAlphaEvolveCodingAgent2025a,huAutomatedDesignAgentic2024}，以及从 rollout 中搜索或学习得到的图结构工作流~\citep{zhangAFlowAutomatingAgentic2024,zhouSymbolicLearningEnables2024a}。

AHE 把完整的 harness 作为一个组合整体来调优，而不是只调优单一的可编辑面，因此跨组件的权衡对优化器而言是可读的。
它同时把人类先验保持在最低限度，将方法论留给优化器从 rollout 中自行发现，而不是由人手工写定。
% The optimization target is long-horizon software-engineering work.
我们将在第~\ref{sec:method}~节描述实现这些选择的基底、轨迹分析与迭代过程。
